% !PRINTED 236
\weberpage{236}
%<<CONT
\weq{16}{\left(\frac{2}{\partial}\right)\frac{u\left(\dfrac{c+\partial\omega}{a}\right)}{u(\omega)^n}}

als Wurzeln einer Modulargleichung, und dies ist die \gesperrt{Jacobische}
Modulargleichung\footnote{Ist $M$ der \gesperrt{Jacobische} Multiplikator, so ist
\[
\frac{(kk')^{\frac{n-1}{3}}}{M} = 4^{\frac{n-1}{3}} P_{11}^2 P_{00}^4.
\]}.

Ist $n$ nicht durch 3 teilbar, so nehmen wir

\weq{17}{c \equiv 0\ (\mathrm{mod}\ 3)}

an, wodurch $\varrho$ den Wert 1 erhält.

In gleicher Weise kann man die Transformation der $\eta$-Funktion
[§\,38, (15), (18), (19)] auf die letzte der Gleichungen (6) anwenden und
erhält:

\weq{18}{P_{11} = \left(\frac{\beta}{\alpha}\right)\left(\frac{\beta'}{\alpha'}\right)\varrho\, i^{\frac{a-1}{2}}\sqrt{\partial}\,\frac{\eta\left(\dfrac{c+\partial\omega}{a}\right)}{\eta(\omega)}}

(wobei es schon genügen würde, wenn $c$ durch 8 teilbar ist). Ist $n$ durch
3 nicht teilbar und $c$ durch 3 teilbar, so ist auch hierin $\varrho=1$ zu
setzen. Die Bestimmung des Vorzeichens in (18) hat für uns nur in dem Falle
Interesse, wo $n$ ein Quadrat ist. Es ist aber [nach (8)]
\[
\left(\frac{\beta}{\alpha}\right)\left(\frac{\beta'}{\alpha'}\right) = \left(\frac{\partial}{\alpha}\right)\left(\frac{\beta'}{\alpha\alpha'}\right) = \left(\frac{\partial}{\alpha}\right)\left(\frac{\beta'}{a}\right) = \left(\frac{\alpha}{\partial}\right)\left(\frac{\beta'}{a}\right)
\]
[letzteres nach dem Reziprozitätsgesetz der quadratischen Reste, weil
$\alpha \equiv 1\ (\mathrm{mod}\ 4)$]. Wenn nun $n$ ein Quadrat ist, so sind
auch $a\!:\!e$, $\partial\!:\!e$ Quadrate und es ergibt sich:
\[
\left(\frac{\alpha}{\partial}\right)\left(\frac{\beta'}{a}\right) = \left(\frac{\alpha\beta'}{e}\right) = \left(\frac{c}{e}\right),
\]
also wird in diesem Falle
\[
P_{11} = \varrho\, i^{\frac{a-1}{2}}\left(\frac{c}{e}\right)\sqrt{\partial}\,\frac{\eta\left(\dfrac{c+\partial\omega}{a}\right)}{\eta(\omega)}.
\]

Die zur Charakterisierung einer Reihe aufgestellten Formeln (10) sind ein
spezieller Fall eines allgemeineren Systems, das man erhält, indem man auf
$\alpha'$, $\beta'$ in (10) eine lineare Transformation anwendet. Man
erhält dann folgendes:
%CONT>>
